% GENERATED FILE. Edit canonical lessons, then run npm run generate:books.
% canonical-source-hash: fnv1a64:00f4b086716314eb
% canonical-lessons: ES-C588-por-tanto, ES-C588-o-sea, ES-C588-en-cambio, ES-C588-sobre-todo, ES-C588-en-general

\chapter{Ordenar las ideas}
\label{ch:es-ordenar-las-ideas}
\hlchaptermodality{588}

\begin{chapteropening}
\textbf{By the end of this chapter:} \emph{I can explain what I mean, contrast two points, stress what matters most and draw a conclusion.}

Complete the last lesson of chapter 588: I can explain what I mean, contrast two points, stress what matters most and draw a conclusion.
\end{chapteropening}

\section[por tanto]{por tanto — therefore, so}

\label{lesson:ES-C588-por-tanto}



\begin{quote}
Before the new one: say the Spanish for the certainty, then the Spanish for apparently.
\end{quote}

\subsection*{You'll want to know: por tanto}
\textbf{por tanto} — \textquotedblleft{}therefore, so\textquotedblright{}.

\begin{quote}
No tengo dinero; por tanto, no voy.
\end{quote}

\begin{quote}
\emph{I have no money; therefore, I am not going.}
\end{quote}

Means the same as \emph{por lo tanto} and \emph{así que}, a little more formal.

\begin{cousinweb}
\emph{por tanto} is built on \emph{tanto}, inherited from Latin \emph{tantus} 'so great, so much' (EtymDB).
\end{cousinweb}

\begin{grammarlens}[title={What you've built}]
One more word for telling, explaining and describing at B1.
\end{grammarlens}

\subsection*{Your turn}
\begin{itemize}
\raggedright
  \item \emph{Say it:} \emph{por tanto}
  \item \emph{Say it:} \emph{por tanto}, once more
  \item \emph{Say it:} a sentence of your own with \emph{por tanto}
  \item \emph{Recall:} say the Spanish for the certainty, then the Spanish for apparently, then say \emph{por tanto} again
\end{itemize}

\begin{tcolorbox}[breakable,colback=teal!4,colframe=teal!35!black,title={Before you move on}]
What does por tanto mean? (\textquotedblleft{}Therefore, so\textquotedblright{}.) Say it once more.
\end{tcolorbox}
\section[o sea]{o sea — I mean, that is}

\label{lesson:ES-C588-o-sea}



\begin{quote}
Before the new one: say the Spanish for apparently, then the Spanish for therefore.
\end{quote}

\subsection*{You'll want to know: o sea}
\textbf{o sea} — \textquotedblleft{}I mean, that is\textquotedblright{}.

\begin{quote}
Llego el lunes, o sea, dentro de tres días.
\end{quote}

\begin{quote}
\emph{I arrive on Monday, that is, in three days.}
\end{quote}

Very common in speech to explain or correct what you said.

\begin{cousinweb}
\emph{o sea} 'or let it be' is built on \emph{o}, inherited from Latin \emph{aut} 'or' (EtymDB), and \emph{sea}, the subjunctive of \emph{ser}, usually traced to Latin \emph{sedeat}, from \emph{sedere} 'to sit'.
\end{cousinweb}

\begin{grammarlens}[title={What you've built}]
One more word for telling, explaining and describing at B1.
\end{grammarlens}

\subsection*{Your turn}
\begin{itemize}
\raggedright
  \item \emph{Say it:} \emph{o sea}
  \item \emph{Say it:} \emph{o sea}, once more
  \item \emph{Say it:} a sentence of your own with \emph{o sea}
  \item \emph{Recall:} say the Spanish for apparently, then the Spanish for therefore, then say \emph{o sea} again
\end{itemize}

\begin{tcolorbox}[breakable,colback=teal!4,colframe=teal!35!black,title={Before you move on}]
What does o sea mean? (\textquotedblleft{}I mean, that is\textquotedblright{}.) Say it once more.
\end{tcolorbox}
\section[en cambio]{en cambio — whereas, by contrast}

\label{lesson:ES-C588-en-cambio}



\begin{quote}
Before the new one: say the Spanish for therefore, then the Spanish for I mean.
\end{quote}

\subsection*{You'll want to know: en cambio}
\textbf{en cambio} — \textquotedblleft{}whereas, by contrast\textquotedblright{}.

\begin{quote}
Yo trabajo; mi hermano, en cambio, estudia.
\end{quote}

\begin{quote}
\emph{I work; my brother, by contrast, studies.}
\end{quote}

It shows a contrast between two people or things.

\begin{cousinweb}
\emph{en cambio} is built on \emph{cambio}, from \emph{cambiar}, which is inherited from Late Latin \emph{cambiare} 'to exchange', a word of Gaulish origin (Lewis \& Short).
\end{cousinweb}

\begin{grammarlens}[title={What you've built}]
One more word for telling, explaining and describing at B1.
\end{grammarlens}

\subsection*{Your turn}
\begin{itemize}
\raggedright
  \item \emph{Say it:} \emph{en cambio}
  \item \emph{Say it:} \emph{en cambio}, once more
  \item \emph{Say it:} a sentence of your own with \emph{en cambio}
  \item \emph{Recall:} say the Spanish for therefore, then the Spanish for I mean, then say \emph{en cambio} again
\end{itemize}

\begin{tcolorbox}[breakable,colback=teal!4,colframe=teal!35!black,title={Before you move on}]
What does en cambio mean? (\textquotedblleft{}Whereas, by contrast\textquotedblright{}.) Say it once more.
\end{tcolorbox}
\section[sobre todo]{sobre todo — above all, especially}

\label{lesson:ES-C588-sobre-todo}



\begin{quote}
Before the new one: say the Spanish for I mean, then the Spanish for whereas.
\end{quote}

\subsection*{You'll want to know: sobre todo}
\textbf{sobre todo} — \textquotedblleft{}above all, especially\textquotedblright{}.

\begin{quote}
Me gusta, sobre todo, el precio.
\end{quote}

\begin{quote}
\emph{I like, above all, the price.}
\end{quote}

Written as two words; \emph{el sobretodo} is an overcoat.

\begin{cousinweb}
\emph{sobre todo} is built on \emph{sobre}, inherited from Latin \emph{super} 'over', and \emph{todo}, from \emph{totus} 'whole' (EtymDB).
\end{cousinweb}

\begin{grammarlens}[title={What you've built}]
One more word for telling, explaining and describing at B1.
\end{grammarlens}

\subsection*{Your turn}
\begin{itemize}
\raggedright
  \item \emph{Say it:} \emph{sobre todo}
  \item \emph{Say it:} \emph{sobre todo}, once more
  \item \emph{Say it:} a sentence of your own with \emph{sobre todo}
  \item \emph{Recall:} say the Spanish for I mean, then the Spanish for whereas, then say \emph{sobre todo} again
\end{itemize}

\begin{tcolorbox}[breakable,colback=teal!4,colframe=teal!35!black,title={Before you move on}]
What does sobre todo mean? (\textquotedblleft{}Above all, especially\textquotedblright{}.) Say it once more.
\end{tcolorbox}
\section[en general]{en general — in general, generally}

\label{lesson:ES-C588-en-general}



\begin{quote}
Before the new one: say the Spanish for whereas, then the Spanish for above all.
\end{quote}

\subsection*{You'll want to know: en general}
\textbf{en general} — \textquotedblleft{}in general, generally\textquotedblright{}.

\begin{quote}
En general, la gente está contenta.
\end{quote}

\begin{quote}
\emph{In general, people are happy.}
\end{quote}

\emph{Por lo general} means the same.

\begin{cousinweb}
\emph{en general} is built on \emph{general}, a learned borrowing from Latin \emph{generalis} 'of a whole class', built on \emph{genus} 'kind' (EtymDB).
\end{cousinweb}

\begin{grammarlens}[title={What you've built}]
One more word for telling, explaining and describing at B1.
\end{grammarlens}

\subsection*{Your turn}
\begin{itemize}
\raggedright
  \item \emph{Say it:} \emph{en general}
  \item \emph{Say it:} \emph{en general}, once more
  \item \emph{Say it:} a sentence of your own with \emph{en general}
  \item \emph{Recall:} say the Spanish for whereas, then the Spanish for above all, then say \emph{en general} again
\end{itemize}

\begin{tcolorbox}[breakable,colback=teal!4,colframe=teal!35!black,title={Before you move on}]
What does en general mean? (\textquotedblleft{}In general, generally\textquotedblright{}.) Say it once more.
\end{tcolorbox}
